%% Chapter2_LiteratureReview.tex
\chapter{Literature Review}
\label{chap:lit}

\section{Historical Overview}
Fixed-time signal plans are commonly derived from Webster's delay-minimizing cycle-length formula~\cite{ref2}. Networks of signals were later coordinated with traffic-responsive systems: SCATS adjusts splits, cycle lengths and offsets from detector data~\cite{ref3}, and SCOOT optimizes the plan on line from loop-detector flows~\cite{ref4}. Vehicle-actuated control, standard in signal timing practice, extends a green while detectors near the stop line keep registering vehicles and ends it when a gap appears (``gap-out''), within minimum and maximum green limits~\cite{ref23}. Max-pressure control chooses phases from queue differences and has provable throughput guarantees~\cite{ref11}.

On the simulation side, the Gipps car-following model~\cite{ref21} and the Krau{\ss} model with its driver reaction time~\cite{ref25} are the basis of many microscopic simulators, SUMO among them~\cite{ref22}. The Highway Capacity Manual~\cite{ref24} supplies the field saturation-flow and start-up lost-time values against which a simulator can be checked. For sensing, Viola and Jones introduced boosted cascades of Haar-like features~\cite{ref16}, which remain attractive for embedded use because inference is cheap and needs no GPU; multi-object trackers combine a Kalman filter~\cite{ref18} with Hungarian assignment~\cite{ref19}, as in ByteTrack~\cite{ref17}.

\section{Recent and State-of-the-Art Work}
Reinforcement-learning controllers have been studied widely, although they must be trained in simulation and their behaviour is harder to bound~\cite{ref12}. Among low-cost density-based systems, Nsofwa and Tembo classify image-based density into low, medium and high levels and adjust signal duration, showing in simulation that this outperforms fixed timing~\cite{ref1}. Kamble et al.\ apply YOLO-based recognition~\cite{ref7} to IP-camera feeds to prioritize the busiest lane~\cite{ref6}; such detectors are accurate but demanding on low-power hardware. Pardeshi et al.\ use passive infrared sensors on an ESP32 platform~\cite{ref5}, and Patil et al.\ and Thomas et al.\ also adapt timing to measured density~\cite{ref9,ref10}. Hu et al.\ estimate density from surveillance cameras~\cite{ref8}. Related systems add incident detection and alerting~\cite{ref13}, AI-based light design~\cite{ref14} or machine-learning traffic prediction on the Edge Impulse platform~\cite{ref15}. Emergency vehicles have been given signal priority based on visual sensing~\cite{ref20}.

These works mainly report an improvement over fixed timing on a prototype or a few scenarios. This project differs in three ways: it compares against strong baselines on identical traffic, it tests demand up to 1.5 times capacity, and it verifies change-interval safety at every simulation step.

\section{Comparative Analysis of Existing Approaches}

\begin{table}[H]
\centering
\caption{Comparison of Existing Approaches}
\label{tab:litreview}
\small
\begin{tabular}{@{}p{3.3cm}p{1.1cm}p{3.6cm}p{3.1cm}p{3.6cm}@{}}
\toprule
\textbf{Author(s)} & \textbf{Year} & \textbf{Approach} & \textbf{Setting} & \textbf{Limitation (for this study)} \\
\midrule
Webster~\cite{ref2} & 1958 & Delay-optimal fixed cycle and splits & Fixed plan from mean demand & Cannot react to fluctuating or uneven demand \\
Sims and Dobinson~\cite{ref3}; Hunt et al.~\cite{ref4} & 1980, 1982 & SCATS and SCOOT: detector-based adaptive plans & Networks with loop detectors & Needs dedicated detector infrastructure \\
Koonce et al.~\cite{ref23} & 2008 & Vehicle-actuated control with gap-out & Stop-line detectors & Needs detectors; no density score \\
Varaiya~\cite{ref11} & 2013 & Max-pressure phase selection & Queue differences & Moving vehicles also leave the pressure count as a queue discharges \\
Nsofwa and Tembo~\cite{ref1} & 2025 & Three-level image-based density classes set the green & Simulation & Compared with fixed timing; no capacity or safety-interval study reported \\
Pardeshi et al.~\cite{ref5} & 2026 & Infrared sensing and density-based prioritization on ESP32 & Prototype platform & Gain reported against fixed timing \\
Kamble et al.~\cite{ref6} & 2025 & YOLO on IP-camera feeds; busiest lane first & Camera feeds & Detector is demanding on low-power hardware \\
Hu et al.~\cite{ref8} & 2023 & Density estimation from surveillance cameras & Existing camera network & Estimates density; timing logic not the focus \\
Patil et al.~\cite{ref9}; Thomas et al.~\cite{ref10} & 2024, 2025 & Timing adapted to measured density & Prototype or few scenarios & Gain reported against fixed timing \\
\bottomrule
\end{tabular}
\end{table}

\section{Identification of Research Gap}
Based on the review above, the following gaps remain in the existing literature:
\begin{enumerate}
  \item Gap 1: Low-cost density-based controllers are usually compared only with fixed timing, so it is unknown how much of their gain comes from the density logic rather than from replacing an untuned plan.
  \item Gap 2: Behaviour near and above capacity (queue growth, spillback, throughput) and the safety of change intervals under variable greens and emergency preemption are rarely examined.
  \item Gap 3: The gap between idealized sensing and the accuracy of a real, inexpensive detector is seldom quantified, nor are the density weights validated.
\end{enumerate}
The proposed project directly addresses these gaps by comparing against re-implemented Webster, actuated, max-pressure and density-level strategies on identical traffic, by sweeping demand from 0.6 to 1.5 times capacity with a per-step safety audit, and by measuring the effect of imperfect sensing and of calibrating the weights on held-out seeds.
